\documentclass[conference]{IEEEtran}
\IEEEoverridecommandlockouts
% Template version as of 6/27/2024

\usepackage{cite}
\usepackage{amsmath,amssymb,amsfonts}
\usepackage{algorithmic}
\usepackage{graphicx}
\usepackage{textcomp}
\usepackage{xcolor}
\usepackage{url}
\usepackage{booktabs}
\def\BibTeX{{\rm B\kern-.05em{\sc i\kern-.025em b}\kern-.08em
    T\kern-.1667em\lower.7ex\hbox{E}\kern-.125emX}}
\begin{document}

\title{Semi-Supervised Learning for Skin Lesion Classification:
Comparing CNNs to Traditional ML Baselines\\
\large CS~4210 Final Project Report}

\author{\IEEEauthorblockN{Andrew Mazmanian, Abigail Moran, Eric Pham,
Khoi Pham, and Sofia Truong}
\IEEEauthorblockA{\textit{Department of Computer Science} \\
\textit{California State Polytechnic University, Pomona}\\
Pomona, USA \\
\{amazmanian, abigailmoran, ericpham, khoinpham, satruong\}@cpp.edu}}

\maketitle

\begin{abstract}
Detecting skin cancer with machine learning usually requires large
amounts of expertly labeled images, which are expensive and difficult
to obtain in clinical practice. In this project we use the HAM10000
dermatoscopic dataset and artificially withhold most labels in the
training partition to simulate a real-world setting where data is
mostly unlabeled. We build a self-training Convolutional Neural
Network (CNN) pipeline and compare it against two traditional
baselines (K-Nearest Neighbors and Decision Tree) and a fully
supervised CNN baseline. Our data-preparation pipeline produces a
stratified train/validation/test split of $5{,}228 / 1{,}121 / 1{,}121$
samples after lesion-aware deduplication, with a configurable
labeled/unlabeled partition of 520 labeled and 4{,}708 unlabeled
training samples. On the held-out test set, KNN obtains 0.733
accuracy and 0.20 macro-F1, the Decision Tree obtains 0.669 accuracy
and 0.21 macro-F1, the ResNet-18 supervised baseline obtains 0.712
accuracy and 0.24 macro-F1, and the ResNet-18 with iterative
self-training obtains 0.733 accuracy and 0.16 macro-F1. Contrary to
our working hypothesis, self-training degrades macro-F1 relative to
the supervised CNN baseline. Under heavy class imbalance the
iterative pseudo-labeling loop reinforces the model's bias toward the
majority class \emph{nv}: the final model predicts only \emph{nv}
(96\%) and \emph{bkl} (4\%) on the test set, and the modest
minority-class recall the supervised baseline had achieved on
\emph{bkl} (0.41) collapses to 0.20. The result is a textbook instance
of the confirmation-bias failure mode of self-training and motivates
class-aware pseudo-label selection as the natural next step.
\end{abstract}

\begin{IEEEkeywords}
Convolutional Neural Networks, Image Classification, Machine Learning,
Self-Training, Semi-Supervised Learning, Skin Cancer.
\end{IEEEkeywords}

\section{Introduction}
Skin cancer is one of the most prevalent cancers worldwide, but it is
highly treatable when detected early, even for the more
life-threatening variants such as melanoma. The problem lies in the
necessity for early detection. Current practice involves physicians
manually evaluating lesion images using dermoscopic imaging, and
considering asymmetry, border irregularity, color variation, diameter,
and evolution by hand. Even with support from modern technology, this
can be time-consuming, subject to the practitioner's experience, and
error-prone due to the variable and complex nature of skin lesions.

The use of classification models as a diagnostic tool is actively
being studied, and although considerable improvements have been made,
limitations still exist within traditional baseline models that use
the K-Nearest Neighbor (KNN) and Decision Tree algorithms. While
these methods can perform reasonably well when paired with carefully
engineered features, they do not inherently learn the complex visual
patterns that exist in skin lesion images, which limits their
effectiveness \cite{wu2022skin}.

Convolutional Neural Networks (CNNs) are considered more effective
for this task, since they automatically learn hierarchical visual
features directly from raw image data without manual feature
engineering. Their effectiveness comes with a practical constraint:
they require large volumes of labeled training data, which can be
difficult to obtain in a clinical setting. Another challenge is the
class imbalance present in dermatological datasets. Semi-supervised
approaches such as self-training can address these issues by using
unlabeled data alongside a small labeled set, gradually expanding
the training pool as model confidence improves \cite{wu2022skin}.

\textbf{Challenges.} Three challenges drive the design of this project:
(1) severe class imbalance in HAM10000, where melanocytic nevi (\texttt{nv})
account for over 60\% of samples; (2) data leakage risk, since multiple
images may share the same physical \texttt{lesion\_id}; and (3) limited
labeled data, simulated in our setup by retaining only a small fraction
of the training labels.

\textbf{Road map.} The remainder of this report is organized as follows.
Section~\ref{sec:related} surveys related work and reported baselines
for HAM10000. Section~\ref{sec:dataset} describes the dataset,
preprocessing, and the resulting splits. Section~\ref{sec:method}
details our four model conditions and the self-training procedure.
Section~\ref{sec:results} reports our completed data-preparation
results and outlines the experimental protocol for model evaluation.
Section~\ref{sec:conclusion} summarizes achievements and future
directions, and Section~\ref{sec:supp} provides links to the LaTeX
source and code repository.

\section{Related Work}
\label{sec:related}
Wu \textit{et al.}~\cite{wu2022skin} provide a systematic review of
deep-learning approaches to skin cancer classification, reporting
that CNN-based methods consistently outperform classical
feature-engineering pipelines (KNN, SVM, decision trees) on
dermatoscopic datasets, while also highlighting class imbalance and
limited labeled data as recurring obstacles.

The HAM10000 dataset itself was introduced by Tschandl
\textit{et al.}~\cite{tschandl2018,b3} as a large multi-source
collection of dermatoscopic images, and was adopted as the basis
of the ISIC 2018 Skin Lesion Analysis Challenge
\cite{codella2019}. In ISIC 2018 Task~3, top submissions reached
balanced multi-class accuracies in the 70--90\% range on the held-out
test set; the challenge organizers report macro-averaged metrics to
discourage models from exploiting the dominance of \texttt{nv}.

Public Kaggle baselines on HAM10000 \cite{kaggle2018} typically
report ResNet-style CNNs in the 75--85\% top-1 accuracy range under
fully supervised settings, while shallow baselines (KNN, Decision
Tree) on flattened or PCA-reduced pixels report substantially lower
macro-F1 scores, often dominated by the majority class. These
patterns motivate our four-way comparison: two classical baselines,
one supervised CNN, and one self-training CNN.

\section{Dataset}
\label{sec:dataset}

\begin{table}[h]
\centering
\caption{Image Data Characteristics of the HAM10000 Dataset}
\begin{tabular}{ll}
\hline
\textbf{Attribute} & \textbf{Description} \\ \hline
Feature Type & Raw dermatoscopic images \\
Total Samples & 10,015 images \\
Image Content & Single centered skin lesion per image \\
Image Type & RGB dermatoscopic images \\
Original Resolution & Approximately $450 \times 600$ pixels \\
Preprocessed Resolution & $224 \times 224 \times 3$ \\
Format & JPEG (.jpg) \\
\hline
\end{tabular}
\end{table}

The HAM10000 (Human Against Machine with 10{,}000 training images)
dataset \cite{tschandl2018} is a widely used benchmark for skin
lesion classification using machine learning. To address the lack
of large, diverse, annotated dermatoscopic datasets in medical
imaging research, HAM10000 is multi-sourced from different
institutions and imaging modalities. Methods of data procurement,
established as ``ground truths,'' are dominated by histopathology
(\textit{histo}, over 50\%), with the remainder coming from follow-up
examinations (\textit{follow up}), expert consensus
(\textit{consensus}), and in-vivo confocal microscopy
(\textit{confocal}).

The 10{,}015 images are categorized into seven clinically significant
classes of pigmented skin lesions: actinic keratoses /
intraepithelial carcinoma (\texttt{akiec}), basal cell carcinoma
(\texttt{bcc}), benign keratosis-like lesions (\texttt{bkl}),
dermatofibroma (\texttt{df}), melanoma (\texttt{mel}), melanocytic
nevi (\texttt{nv}), and vascular lesions (\texttt{vasc}).

\subsection{Dataset Characteristics and Limitations}
Two notable limitations require careful handling to ensure optimal
model accuracy: severe class imbalance (with melanocytic nevi
comprising over 60\% of samples) and potential data leakage from
multiple images sharing the same \texttt{lesion\_id}
\cite{tschandl2018,b3,kaggle2018}. While the input feature (image
size) varies based on the modeling approach, the target output is
consistently one of the seven diagnostic labels.

\subsection{Preprocessing and Splits}
\label{sec:splits}
Our preparation pipeline performs the following steps:
\begin{enumerate}
    \item Verify that every image referenced in the metadata file
    exists on disk.
    \item Drop duplicate images sharing the same \texttt{lesion\_id},
    keeping a single representative per lesion.
    \item Encode the seven diagnoses into integer labels $\{0,\ldots,6\}$.
    \item Resize all images to $224 \times 224$ pixels and normalize
    using the ImageNet channel-wise mean
    $\mu = [0.485, 0.456, 0.406]$ and standard deviation
    $\sigma = [0.229, 0.224, 0.225]$, consistent with our pretrained
    backbone.
    \item Split the deduplicated set into training (70\%),
    validation (15\%), and test (15\%) using stratified sampling
    on the diagnosis label.
    \item Within the training partition, sample a labeled subset
    per class, treating the rest as unlabeled. The split is
    configurable through \texttt{LABELED\_DATA\_PERCENTAGE} in
    \texttt{config.py} so we can sweep this ratio in experiments.
\end{enumerate}

After deduplication, $2{,}545$ duplicate lesion images were removed,
leaving $7{,}470$ unique lesions. Table~\ref{tab:splits} summarizes
the resulting partitions, and Table~\ref{tab:classdist} reports the
per-class counts.

\begin{table}[h]
\centering
\caption{Data partitions after deduplication.}
\label{tab:splits}
\begin{tabular}{lr}
\hline
\textbf{Partition} & \textbf{\#Samples} \\ \hline
Full (after dedup)         & 7{,}470 \\
Training                   & 5{,}228 \\
Validation                 & 1{,}121 \\
Test                       & 1{,}121 \\
\quad Labeled training     & 520 \\
\quad Unlabeled training   & 4{,}708 \\
\hline
\end{tabular}
\end{table}

\begin{table}[h]
\centering
\caption{Per-class counts after lesion-level deduplication.}
\label{tab:classdist}
\begin{tabular}{lr}
\hline
\textbf{Class} & \textbf{\#Samples} \\ \hline
\texttt{nv}    & 5{,}403 \\
\texttt{bkl}   &   727 \\
\texttt{mel}   &   614 \\
\texttt{bcc}   &   327 \\
\texttt{akiec} &   228 \\
\texttt{vasc}  &    98 \\
\texttt{df}    &    73 \\
\hline
\textbf{Total} & \textbf{7{,}470} \\
\hline
\end{tabular}
\end{table}

The final partitions clearly preserve the long-tailed class
distribution: \texttt{nv} dominates while \texttt{vasc} and
\texttt{df} are heavily under-represented. This motivates the use of
class-weighted losses, stratified sampling, and macro-averaged
metrics throughout our experiments.

\section{Methodology}
\label{sec:method}

\subsection{Dataset Preparation and Preprocessing}
The pipeline described in Section~\ref{sec:splits} is implemented
in \texttt{data\_preparation.py} and \texttt{utils.py}. To make the
project portable across collaborators' machines, the saved CSVs
store image paths relative to the project root and the loader
resolves them at runtime, so no machine-specific absolute paths
appear in the repository.

\subsection{Train/Validation/Test Split and Semi-Supervised Setup}
The dataset is split into training (70\%), validation (15\%), and
test (15\%) using stratified sampling. To simulate a semi-supervised
setting, only a small fraction of the training labels is retained,
while the remainder is treated as unlabeled \cite{codella2019}. The
validation and test sets remain fully labeled. This setup enables
controlled comparison of supervised baselines and self-training CNN
models under limited-label conditions.

\subsection{Machine Learning Models}

\subsubsection{Traditional Supervised Baselines}
To establish benchmarks under labeled-only conditions, we train two
classical classifiers on the labeled split. Each image is flattened
and reduced to 100 components via Principal Component Analysis
(PCA) before classification.

\textit{K-Nearest Neighbors (KNN):} Trained on PCA-reduced features
using Euclidean distance, with $k$ selected via 5-fold cross-validated
grid search on the training set over $k \in \{3, 5, 7, 11\}$.

\textit{Decision Tree:} Trained using the Gini impurity criterion,
with maximum depth tuned over $\{5, 10, 20, \text{None}\}$ on the
validation set.

These baselines illustrate the performance ceiling of traditional
methods on complex image data under limited-label conditions.

\subsubsection{CNN Supervised Baseline}
We fine-tune a ResNet-18 backbone pretrained on ImageNet, replacing
the final layer with a seven-class output head. The model is
trained on the labeled split using a class-weighted cross-entropy
loss, the Adam optimizer with learning rate $1 \times 10^{-4}$,
and a cosine annealing scheduler over 20 epochs. Data augmentation
includes random horizontal flips, rotation up to $20^{\circ}$, and
color jitter. The implementation in
\texttt{train\_supervised\_baseline.py} follows this recipe with one
pragmatic simplification we report transparently: the class weights
are set to a fixed vector $[0.2, 1, 1, 1, 1, 1, 1]$ that downweights
the majority class \texttt{nv} rather than using exact
inverse-frequency weights.

\subsubsection{CNN with Self-Training}
Our primary model extends the CNN baseline by incorporating
unlabeled data through iterative self-training:
\begin{enumerate}
    \item Train the CNN on the labeled set.
    \item Generate pseudo-labels for all unlabeled samples using
    the model's softmax predictions.
    \item Retain samples where the maximum predicted probability
    exceeds threshold $\tau$ (initialized at $\tau = 0.90$).
    \item Retrain the CNN on the combined labeled and pseudo-labeled set.
    \item Reduce $\tau$ by 0.05 and repeat for three total iterations.
\end{enumerate}
This confidence-threshold decay allows the model to incorporate
pseudo-labeled data conservatively at first, gradually admitting
more unlabeled samples as predictions improve.

\subsection{Model Evaluation}
All models are evaluated on the held-out test set. Following the
ISIC 2018 evaluation strategy \cite{codella2019}, macro-averaged
F1 score is the primary metric to weight all seven classes equally
regardless of support. We additionally report overall accuracy,
per-class one-vs-rest AUC-ROC, and confusion matrices to capture
systematic misclassification patterns. We compare four conditions:
KNN, Decision Tree, CNN supervised baseline, and CNN with
self-training, isolating the contributions of the deep-learning
architecture and the semi-supervised procedure to overall
classification performance.

To guarantee that all four conditions are scored identically, we
implemented a shared utility \texttt{eval\_script.py} exposing a
single function \texttt{evaluate(model\_name, y\_true, y\_pred, y\_prob)}.
Given the test labels, predicted classes, and predicted class
probabilities from any model, it computes macro-F1, accuracy, the
seven per-class AUC values, and the confusion matrix, and saves a
combined figure (confusion matrix and one-vs-rest ROC curves) to
\texttt{<model\_name>\_eval.png}. Each modeling script imports this
function so the report tables and figures are produced from a
single, version-controlled implementation.

\section{Results}
\label{sec:results}

\subsection{Data-Preparation Results}
The data-preparation pipeline executes end-to-end on a standard
laptop in well under a minute. Of the $10{,}015$ rows in
\texttt{HAM10000\_metadata}, all images were located on disk
(0 missing). Lesion-level deduplication removed $2{,}545$ duplicate
images, yielding $7{,}470$ unique lesions distributed as in
Table~\ref{tab:classdist}. Stratified splitting produced the
$5{,}228 / 1{,}121 / 1{,}121$ partition reported in
Table~\ref{tab:splits}. The labeled/unlabeled split inside the
training partition produced $520$ labeled and $4{,}708$ unlabeled
images at the current configuration, with at least
\texttt{MIN\_SAMPLES\_PER\_CLASS} examples retained for every class
to prevent any class from being entirely unlabeled.

\subsection{CNN Supervised Baseline: Training Progress}
The CNN supervised baseline has been implemented and trained for the
full 20 epochs on the 520-sample labeled training set. Each epoch
iterates over the labeled loader with batch size 32, applies the
augmentation pipeline (flip, rotation, color jitter), computes
class-weighted cross-entropy loss, and updates the weights with
Adam ($\mathrm{lr} = 1 \times 10^{-4}$). The training loop logs the
per-epoch mean loss and saves the final weights to
\texttt{resnet18\_supervised\_baseline.pth} (approximately 44.8\,MB).
This checkpoint also serves as the warm-start for the self-training
iterations described in Section~IV-C3, and its full test-set metrics
are reported in Table~\ref{tab:results} alongside the three other
conditions.

\subsection{Experimental Protocol for Models}
With the splits above frozen, the models are trained on the 
labeled training set and evaluated once on the test set. 
With the exception of the KNN baseline (which utilizes internal 
5-fold cross-validation on the training set), hyperparameters for the 
remaining conditions are explicitly tuned on the validation set. For the self-training
condition we additionally use the unlabeled training set during the
iterative pseudo-labeling loop, warm-started from the supervised
baseline checkpoint described above. Macro-F1, accuracy, per-class
AUC-ROC, and confusion matrices are reported in Table~\ref{tab:results}.

\subsection{Comparative Test-Set Results}
All four planned conditions have now been evaluated on the held-out
test set (1{,}121 samples) using the shared \texttt{eval\_script.py}
utility. Macro-F1, accuracy, and per-class one-vs-rest AUC are
computed from the same confusion matrices and probability outputs;
the corresponding figures are saved as
\texttt{KNN\_eval.png},
\texttt{Decision\_Tree\_Baseline\_eval.png},
\texttt{ResNet-18 Supervised Baseline\_eval.png}, and
\texttt{ResNet-18 Self-Training\_eval.png} in the repository.
Table~\ref{tab:results} summarises the headline metrics.

\begin{table}[h]
\centering
\caption{Test-set results (1{,}121 samples). Per-class AUC is the
mean of the seven one-vs-rest values. The self-training row reports
the final model after three iterations with $\tau \in \{0.90, 0.85,
0.80\}$.}
\label{tab:results}
\begin{tabular}{lccc}
\hline
\textbf{Model} & \textbf{Macro-F1} & \textbf{Accuracy} \\
\hline
KNN (PCA + Euclidean)         & 0.2044 & 0.7333 \\
Decision Tree (PCA + Gini)    & 0.2103 & 0.6690 \\
ResNet-18 supervised baseline & 0.2357 & 0.7119 \\
ResNet-18 self-training       & 0.1624 & 0.7333 \\
\hline
\end{tabular}
\end{table}

\subsection{Findings}
The self-training condition disconfirms our working hypothesis. With
$\tau$ initialised at 0.90 and decayed by 0.05 across three
iterations, the model warm-started from the supervised baseline
checkpoint and was retrained on the union of the 520 labeled samples
and every unlabeled sample whose top softmax probability exceeded the
current threshold. Despite this conservative schedule, macro-F1 fell
from 0.2357 to 0.1624, while accuracy rose marginally from 71.2\% to
73.3\%. The accuracy gain is illusory: the final model predicts
\emph{nv} for 1{,}078 of the 1{,}121 test samples (96.2\%) and
\emph{bkl} for the remaining 43, and produces \emph{zero} predictions
for \emph{mel}, \emph{bcc}, \emph{akiec}, \emph{vasc}, and \emph{df}.
The recall on \emph{bkl}, which the supervised baseline had pushed
from 0.19 to 0.41, collapses back to 0.20.

Two observations explain this collapse. First, the supervised baseline
was already biased toward \emph{nv} (79\% of predictions); when used
to pseudo-label the 4{,}708 unlabeled images, it therefore generated
a pseudo-label pool that itself was dominated by \emph{nv}, so each
self-training iteration ingested progressively more \emph{nv}
examples and reinforced the bias rather than correcting it. Second,
the high-confidence threshold ($\tau \geq 0.80$) selects exactly the
samples on which the model is already most certain --- which under
class imbalance are disproportionately the easy majority-class
samples. The mean per-class AUC on the test set drops from 0.66
(supervised) to 0.51 (self-training), confirming that the post-loop
probability ranking is barely better than chance for the minority
classes. This is the textbook confirmation-bias failure mode of
self-training under heavy class imbalance, and it is precisely the failure mode our methodology section flagged as the central empirical risk.

\subsection{Final Ranking and Failure Analysis}
The completed ranking on macro-F1 is therefore
\textbf{ResNet-18 self-training (0.16) $<$ KNN (0.20) $<$ Decision Tree
(0.21) $<$ ResNet-18 supervised (0.24)}, which inverts the ordering
predicted in Section~II for the top two methods. The expected ordering
held among the three classical or single-pass conditions: the supervised
CNN beat both PCA-based baselines by a small but consistent margin on
both macro-F1 and minority-class recall. What did not hold is the
assumption that adding 4{,}708 unlabeled images via self-training would
help; under our current configuration it harms the model. Two design
choices, in retrospect, made this almost guaranteed: (a) pseudo-labels
were generated by a model that was itself heavily majority-biased, and
(b) the confidence-threshold rule selects easy (majority) samples first,
so the loop has no mechanism to inject minority-class signal. We discuss
class-aware fixes in Section~VI.

\section{Conclusion}
\label{sec:conclusion}
\textbf{Achievements.} We implemented and validated a reproducible
data-preparation pipeline for HAM10000 that handles dataset
verification, lesion-aware deduplication, label encoding, stratified
train/validation/test splitting, and a configurable semi-supervised
labeled/unlabeled split. The pipeline is portable across machines
(relative paths in CSVs), deterministic (fixed random seed), and
parameterised through a single \texttt{config.py} file. All four
planned models have been trained and evaluated end-to-end on the
held-out test set using a shared evaluation utility: a KNN baseline
on PCA-reduced features (0.20 macro-F1, 73.3\% accuracy), a Decision
Tree baseline (0.21 macro-F1, 66.9\% accuracy), a ResNet-18
supervised CNN baseline (0.24 macro-F1, 71.2\% accuracy), and a
ResNet-18 with three rounds of confidence-decayed self-training
(0.16 macro-F1, 73.3\% accuracy). The most consequential empirical
finding is negative: under our configuration, self-training
\emph{worsens} macro-F1 by amplifying the supervised model's
existing bias toward the majority class \emph{nv} via repeated
high-confidence pseudo-labeling.

\textbf{Future directions.} The self-training collapse points to
three concrete fixes that we expect would reverse the result: (i)
\emph{class-balanced pseudo-label selection}, e.g.\ keeping the
top-$k$ most confident pseudo-labels per predicted class rather than
a global threshold, so minority classes are forced into the training
pool; (ii) \emph{class-aware thresholds} $\tau_c$ tuned per class, so
that minority-class pseudo-labels are admitted at lower confidence
than majority-class ones, as proposed in FlexMatch-style methods;
and (iii) stronger \emph{class-rebalancing} during retraining (e.g.\
exact inverse-frequency weights or oversampling of pseudo-labeled
minority samples). Beyond these, we plan sensitivity analyses on the
labeled-data percentage, the threshold schedule, and the choice of
backbone (ResNet-50, EfficientNet-B0), as well as a qualitative
analysis of the confusion matrices to identify which lesion types
most benefit from unlabeled data once the bias-amplification problem
is fixed.

\section{Supplementary Materials}
\label{sec:supp}
The LaTeX source of this report and the full project code
(\texttt{config.py}, \texttt{utils.py}, \texttt{data\_preparation.py},
\texttt{example\_usage.py},
\texttt{train\_supervised\_baseline.py},
\texttt{evaluate\_supervised\_baseline.py},
\texttt{train\_self\_training.py},
\texttt{evaluate\_self\_training.py},
\texttt{train\_decision\_tree.py},
\texttt{knn.py},
the shared evaluation utility \texttt{eval\_script.py},
the trained checkpoints \texttt{resnet18\_supervised\_baseline.pth}
and \texttt{resnet18\_self\_training.pth},
the prepared CSV splits, the per-model evaluation figures, and
dataset statistics) are available at:
\begin{center}
\url{https://github.com/khoipc305/ML_AI_project}
\end{center}
The repository \texttt{README.md} documents how to download the
HAM10000 images from Kaggle \cite{kaggle2018} and reproduce all
splits with a single command.

\begin{thebibliography}{9}

\bibitem{wu2022skin}
Y. Wu \textit{et al.}, ``Skin Cancer Classification With Deep
Learning: A Systematic Review,'' \textit{Frontiers in Oncology},
vol.~12, 2022, doi: 10.3389/fonc.2022.893972.

\bibitem{tschandl2018}
P. Tschandl, C. Rosendahl, and H. Kittler, ``The HAM10000 dataset,
a large collection of multi-source dermatoscopic images of common
pigmented skin lesions,'' \textit{Sci. Data}, vol.~5, p.~180161,
Aug.~2018, doi: 10.1038/sdata.2018.161.

\bibitem{b3}
P. Tschandl, ``The HAM10000 dataset, a large collection of
multi-source dermatoscopic images of common pigmented skin
lesions,'' \textit{Harvard Dataverse}, vol.~4, 2018,
doi: 10.7910/DVN/DBW86T.

\bibitem{kaggle2018}
K. S. Mader, ``Skin Cancer MNIST: HAM10000,'' Kaggle, 2018.
[Online]. Available:
\url{https://www.kaggle.com/datasets/kmader/skin-cancer-mnist-ham10000}

\bibitem{codella2019}
N. Codella \textit{et al.}, ``Skin lesion analysis toward melanoma
detection 2018: A challenge hosted by the International Skin Imaging
Collaboration (ISIC),'' arXiv:1902.03368, 2019. [Online]. Available:
\url{https://arxiv.org/abs/1902.03368}

\bibitem{he2016resnet}
K. He, X. Zhang, S. Ren, and J. Sun, ``Deep Residual Learning for
Image Recognition,'' in \textit{Proc. IEEE Conf. Computer Vision
and Pattern Recognition (CVPR)}, 2016, pp.~770--778,
doi: 10.1109/CVPR.2016.90.

\bibitem{lee2013pseudo}
D.-H. Lee, ``Pseudo-Label: The Simple and Efficient Semi-Supervised
Learning Method for Deep Neural Networks,'' in \textit{ICML
Workshop on Challenges in Representation Learning}, 2013.

\bibitem{deng2009imagenet}
J. Deng \textit{et al.}, ``ImageNet: A Large-Scale Hierarchical
Image Database,'' in \textit{Proc. IEEE Conf. Computer Vision and
Pattern Recognition (CVPR)}, 2009, pp.~248--255,
doi: 10.1109/CVPR.2009.5206848.

\end{thebibliography}

\end{document}
